\chapter{Le problème : ce que subissent assureurs et promoteurs}
\label{ch:douleurs}

Ce chapitre décrit les douleurs telles qu'elles se manifestent dans le travail
réel, formulées en \textbf{coûts}, en \textbf{délais} et en \textbf{risques
juridiques}. Chacune est ensuite reprise en \S\ref{ch:reponse} sous la forme
« douleur \(\rightarrow\) réponse technique ».

\section{Sept douleurs}

\subsection{D1 — L'adresse n'est pas l'unité des données publiques}

Les aléas sont publiés à la commune ou par emprises cartographiques ; le
bâtiment, lui, est décrit par la BDNB. Répondre « au bâtiment » avec une donnée
communale produit une réponse qui paraît précise et ne l'est pas.

\begin{douleur}[title={D1 — Granularité inadaptée}]
\textbf{Effet} : sur-estimation sur les biens situés hors des emprises, ou
sous-estimation sur ceux situés à l'intérieur d'une emprise communale large.
\textbf{Coût} : reprise manuelle pour les cas litigieux, et perte de confiance
dans l'outil dès le premier écart détecté par un expert.
\end{douleur}

\subsection{D2 — Aucune clé de jointure commune}

Le géocodeur rend une adresse et un point ; la BDNB rend un
\code{batiment\_groupe\_id} et une géométrie ; Géorisques rend des polygones
d'aléas ; l'hydrographie rend des tronçons ; la photo rend une prise de vue. Il
n'existe pas d'identifiant partagé : la jointure se fait par
\textbf{géométrie} (point dans polygone) et par codes administratifs, avec toutes
les approximations que cela implique.

\subsection{D3 — Résolution hétérogène, non déclarée}

Certaines données sont communales (l'aléa retrait-gonflement des argiles), la
plupart sont vectorielles et ponctuelles, quelques-unes sont surfaciques. Mélanger
ces niveaux sans le dire revient à présenter une estimation communale comme une
mesure sur le bien.

\begin{invariant}[title={Chaque fait déclare sa finesse}]
Tout objet d'aléa du contrat porte \code{source}, \code{resolution} et
l'attribution LO 2.0 du producteur. L'aléa argile est \emph{toujours}
\code{commune-level-estimate} tant qu'aucune source vectorielle n'est câblée :
la changer exige un amendement de la constitution du projet.
\end{invariant}

\subsection{D4 — « Absence de donnée » confondue avec « absence de risque »}

C'est la douleur la plus coûteuse, parce qu'elle est \emph{invisible}. Une source
qui répond 500, un quota atteint, un service momentanément indisponible : si
l'outil rend alors « aucun aléa », l'utilisateur conclut à un bien sans risque.
Le pire des deux mondes — une décision prise sur du vide, avec l'apparence d'une
réponse.

\subsection{D5 — Traçabilité insuffisante pour justifier une décision}

Un assureur doit pouvoir expliquer une tarification. Une valeur sans producteur,
sans date de millésime et sans méthode n'est pas défendable devant un régulateur
ni devant un client qui conteste.

\subsection{D6 — Le coût et le délai de l'analyse manuelle}

Croiser six sources par bien, à la main, représente typiquement plusieurs heures
par dossier. Cela rend l'analyse exhaustive d'un portefeuille économiquement
impossible : on échantillonne, donc on subit le risque de queue.

\subsection{D7 — La dette juridique de l'IA générative}

Introduire un modèle de langage dans un document réglementaire sans garde-fou
revient à accepter qu'un chiffre inventé soit publié. La douleur n'est pas
technique, elle est juridique : la responsabilité du contenu reste à celui qui le
signe.

\section{Les douleurs, mises en regard}

\begin{table}[h]
\centering
\small
\begin{tabularx}{\textwidth}{@{}p{4.2cm}Xp{3.4cm}@{}}
\toprule
\textbf{Douleur} & \textbf{Manifestation observable} & \textbf{Ce qu'elle coûte} \\
\midrule
D1 Granularité & Écarts entre l'outil et l'expert & Crédibilité \\
D2 Jointure & Scripts ad hoc, cas non couverts & Maintenance \\
D3 Résolution & Un chiffre communal présenté comme un fait du bien & Risque juridique \\
D4 Absence \(\neq\) 0 & Décision prise sur une source muette & Sinistre non anticipé \\
D5 Traçabilité & Valeur sans producteur ni millésime & Non-conformité \\
D6 Volume & Analyse limitée à un échantillon & Risque de queue \\
D7 IA non bornée & Chiffre plausible mais faux & Responsabilité \\
\bottomrule
\end{tabularx}
\caption{Les sept douleurs, leur manifestation et ce qu'elles coûtent réellement.}
\end{table}

\section{Le test de réalité : un portefeuille de 1 200 biens}

Pour rendre le problème tangible, prenons un cas représentatif.

Un assureur doit réévaluer un portefeuille de 1 200 biens répartis sur le
territoire. Le travail d'analyse disponible représente l'équivalent de deux
postes à temps plein. À plusieurs heures par bien, seule une fraction du
portefeuille peut être instruite en profondeur ; le reste est traité par
approximation communale.

Trois conséquences s'enchaînent :

\begin{enumerate}
  \item \textbf{Sur les biens non instruits}, la tarification retenue est celle
    de la commune — par exemple une parcelle en bord de cours d'eau traitée
    comme une parcelle quelconque du fait d'une emprise mal comprise.
  \item \textbf{Sur les biens contestés}, l'analyse doit être refaite à la main
    pour produire une justification acceptable, sans capitalisation du travail
    déjà fait.
  \item \textbf{Sur l'ensemble}, aucune réponse n'est reproductible : deux
    analystes, deux méthodes, deux résultats.
\end{enumerate}

\begin{encadre}[title={Ce que la réponse doit donc garantir}]
Une réponse \emph{par bâtiment}, \emph{traçable}, \emph{reproductible}, et
capable de dire « je ne sais pas » sans le confondre avec « il n'y a rien ».
C'est exactement le cahier des charges de \S\ref{ch:reponse}.
\end{encadre}
